% Generated by scripts/assemble.py; corrections belong in source/corrections.json.
\examhead{139}{华中师范大学 2025 年数学分析试卷}{132}
% p133-b19; source PDF page 133
\par\Needspace{5\baselineskip}
\blockanchor{p133-b19}
\majorquestion{一}计算题. 每题 10 分, 共 50 分.\par
% p133-b20; source PDF page 133
\blockanchor{p133-b20}
\question{1}求极限 $\lim_{R\to +\infty}\iint \limits_{D}(x^{2} + y^{2})e^{-(x^{2} + y^{2})}\mathrm{d}x\mathrm{d}y,$ 其中区域 $D = \{(x,y)\mid |x|\leq R,|y|\leq R\}$.\par
% p133-b21; source PDF page 133
\blockanchor{p133-b21}
\question{2}设 $\varphi(t) = \int_0^1 \ln \sqrt{x^2 + t^2} \, \mathrm{d}x$ ($0 \leq t \leq 1$), 求 $\varphi$ 在 $t = 0$ 处的单侧导数.\par
% p133-b22; source PDF page 133
\blockanchor{p133-b22}
\question{3}计算 $J = \int_{0}^{\pi} \frac{x \sin x}{1 + \cos^{2} x} \, \mathrm{d}x$.\par
% p133-b23; source PDF page 133
\blockanchor{p133-b23}
\question{4}求全微分 $(y + z)\mathrm{d}x + (z + x)\mathrm{d}y + (x + y)\mathrm{d}z$ 的原函数 $u(x,y,z)$.\par
% p133-b24; source PDF page 133
\blockanchor{p133-b24}
\question{5}计算 $\iint_{S} z^{3} \mathrm{d}x \mathrm{~d}y$, 其中 $S$ 是椭球面 $\frac{x^{2}}{a^{2}} + \frac{y^{2}}{b^{2}} + \frac{z^{2}}{c^{2}} = 1 (a, b, c > 0)$ 的上半部分, 并取外侧.\par
% p134-b00; source PDF page 134
\par\Needspace{5\baselineskip}
\blockanchor{p134-b00}
\majorquestion{二}(10 分) 用数列极限的柯西准则证明 $\lim_{n\to \infty}\sin n$ 不存在.\par
% p134-b01; source PDF page 134
\par\Needspace{5\baselineskip}
\par\Needspace{6\baselineskip}
\blockanchor{p134-b01}
\majorquestion{三}（15分）设不含原点的开区域 $D\subset \mathbb{R}^2$ 上的二元可微函数 $F$ 在直角坐标系下可表示为\par
\nopagebreak[4]
% p134-b02; source PDF page 134
\blockanchor{p134-b02}
\displayformula{z = F (x, y) = f (x) g (y), f (x) \neq 0, g (y) \neq 0.}
% p134-b03; source PDF page 134
\par\Needspace{6\baselineskip}
\blockanchor{p134-b03}
它在极坐标系可表示为\par
\nopagebreak[4]
% p134-b04; source PDF page 134
\blockanchor{p134-b04}
\displayformula{z = F (r \cos \theta , r \sin \theta) = h (r).}
% p134-b05; source PDF page 134
\blockanchor{p134-b05}
求二元函数  $z = F(x, y)$  的显式表达式.\par
% p134-b06; source PDF page 134
\par\Needspace{5\baselineskip}
\blockanchor{p134-b06}
\majorquestion{四}（15分）考察级数 $\sum_{n=1}^{\infty}\left(1 + \frac{1}{2} + \cdots + \frac{1}{n}\right)\frac{\sin n}{n}$ 的敛散性.\par
% p134-b07; source PDF page 134
\par\Needspace{5\baselineskip}
\par\Needspace{6\baselineskip}
\blockanchor{p134-b07}
\majorquestion{五}(15 分) 设函数 $f$ 在 $\mathbb{R}$ 上的二阶导数 $f''(x) > 0$, 又设存在点 $x_0 \in \mathbb{R}$, 使得 $f(x_0) < 0$, 且\par
\nopagebreak[4]
% p134-b08; source PDF page 134
\blockanchor{p134-b08}
\displayformula{\lim _ {x \to - \infty} f ^ {\prime} (x) = \alpha <   0, \lim _ {x \to + \infty} f ^ {\prime} (x) = \beta > 0.}
% p134-b09; source PDF page 134
\blockanchor{p134-b09}
证明： $f(x)$ 有且仅有2个零点.\par
% p134-b10; source PDF page 134
\par\Needspace{5\baselineskip}
\par\Needspace{6\baselineskip}
\blockanchor{p134-b10}
\majorquestion{六}(10 分) 设函数 $f$ 在有限区间 $[a, b]$ 上连续, 且对任意的 $[\alpha, \beta] \subset [a, b]$, 有\par
\nopagebreak[4]
% p134-b11; source PDF page 134
\blockanchor{p134-b11}
\displayformula{\left| \int_ {\alpha} ^ {\beta} f (x) \mathrm{d} x \right| \leq M | \beta - \alpha | ^ {s + 1}.}
% p134-b12; source PDF page 134
\blockanchor{p134-b12}
其中 M, s 为正常数. 证明:  $f(x) \equiv 0, x \in [a, b]$ .\par
% p134-b13; source PDF page 134
\par\Needspace{5\baselineskip}
\par\Needspace{6\baselineskip}
\blockanchor{p134-b13}
\majorquestion{七}(15 分) 设函数 $f$ 在任意有界区间上可积, 并且 $\lim_{x \to +\infty} f(x) = \alpha$ ($\alpha$ 为有限数). 证明:\par
\nopagebreak[4]
% p134-b14; source PDF page 134
\blockanchor{p134-b14}
\displayformula{\lim _ {y \rightarrow 0 ^ {+}} y \int_ {0} ^ {+ \infty} e ^ {- y x} f (x) \mathrm{d} x = \alpha .}
% p134-b15; source PDF page 134
\par\Needspace{5\baselineskip}
\blockanchor{p134-b15}
\majorquestion{八}(20 分) 给定级数 $\sum_{n=1}^{\infty} \frac{a_n}{n^x}$, 其中 $a_n$ 均为实数. 证明:\par
% p134-b16; source PDF page 134
\blockanchor{p134-b16}
(1) (8 分) 若 $\sum_{n=1}^{\infty} \frac{a_n}{n^\lambda}$ 收敛, 则对任意的 $x > \lambda$, $\sum_{n=1}^{\infty} \frac{a_n}{n^x}$ 收敛.\par
% p134-b17; source PDF page 134
\blockanchor{p134-b17}
(2) (12 分) 存在 $r \in [-\infty, +\infty]$, 使得 $x > r$ 时, $\sum_{n=1}^{\infty} \frac{a_n}{n^x}$ 收敛; $x < r$ 时, $\sum_{n=1}^{\infty} \frac{a_n}{n^x}$ 发散.\par
